\begin{progprob}
    % Write the problem here.

    Let \python{data} be a list of $n$ unique real numbers. Furthermore,
    suppose that each number in \python{data} is assigned a color -- it is
    either \python{'red'} or \python{'blue'}. Let \python{colors} be a list
    of $n$ strings, such that \python{colors[i]} gives the color of
    \python{data[i]}. In all parts of this problem, you may assume for simplicity
    that there is at least one data point of each color.
    
    \begin{subprobset}

        \begin{subprob}
            To begin, suppose that \textbf{all of the blue points are less than all of the red points}.
            In a file named \python{ep02.py}, write an efficient function called
            \python{learn_theta(data, colors)} which takes in two arguments -- the lists \python{data}
            and \python{colors} as described above -- and returns a single number $\theta$
            such that all of the blue points are $\leq
            \theta$ and all of the red points are $> \theta$, as is depicted in
            the picture below. You may not assume that \python{data} is
            sorted. The time complexity of your algorithm should be optimal.

            \begin{center}
            \begin{tikzpicture}
                \draw[<->] (-1,0) -- (10,0);
                \foreach \x in {0, 1, 1.5, 2, 3} {
                    \node (pt) at (\x, 0) {$\times$};
                    \node[below=.3 of pt] {B} edge[->] (pt);
                }
                \foreach \x in {5, 5.5, 7, 7.25, 9} {
                    \node (pt) at (\x, 0) {$\times$};
                    \node[above=.3 of pt] {R} edge[->] (pt);
                }

                \coordinate (theta) at (4.5, 0);
                \node[below left] at (theta) {$\theta$};
                \draw[dashed, thick] ([yshift=2.5em]theta) edge ([yshift=-2.5em]theta);

            \end{tikzpicture}
            \end{center}        

            
            \begin{soln}
                We can set $\theta$ to be any number between the biggest blue
                point and the smallest red point. In fact, we can set $\theta$
                to be \emph{exactly} the largest blue point, as then all blue
                points are $\leq \theta$ as desired.

                We can do this in linear time by looping through all of the
                data points, keeping track of the largest blue point seen
                so far:

                \inputminted{python}{\thisdir/include-solution/pp01a-enum.py}

                A more ``pythonic'' approach is to ditch the \python{range}
                and use \python{zip}. \python{zip} takes two (or more) lists
                and returns pairs by ``zipping'' the lists together, elementwise:

                \inputminted{python}{\thisdir/include-solution/pp01a.py}

                Of course, we could always use a comprehension. The below
                filters out the red points and returns the max of the blues:

                \inputminted{python}{\thisdir/include-solution/pp01a-max.py}

                You could make the above into a one-liner, but I like the
                fact that the intermediate result is stored in \python{blue_points};
                it makes the code more readable, in my opinion. Avoid trying to show
                off by writing one-liners. A lot of people can cram code into
                one line, but it takes thought and care to write clean, readable code.

                All of these solutions are essentially the same, and run in
                linear time.
            \end{soln}

        \end{subprob}
    \end{subprobset}

        Now suppose that a small number of the red points are smaller
        than some blue points -- that is, there is some overlap, as shown
        below. Assume for simplicity that the largest data point is red.

        \begin{center}
        \begin{tikzpicture}
            \draw[<->] (-1,0) -- (10,0);
            \foreach \x in {0, 1, 1.5, 2, 7} {
                \node (pt) at (\x, 0) {$\times$};
                \node[below=.3 of pt] {B} edge[->] (pt);
            }
            \foreach \x in {3, 5.5, 7.25, 9} {
                \node (pt) at (\x, 0) {$\times$};
                \node[above=.3 of pt] {R} edge[->] (pt);
            }

            \coordinate (theta) at (4.5, 0);
            \node[below left] at (theta) {$\theta$};
            \draw[dashed, thick] ([yshift=2.5em]theta) edge ([yshift=-2.5em]theta);

        \end{tikzpicture}
        \end{center}        

        We wish to find a real number $\theta$ which ``best'' separates the
        blue points and red points. Clearly the points cannot be separated
        perfectly. Instead, we define a loss function $L(\theta)$ which
        counts the number of points which are on the wrong side of
        $\theta$.  More precisely:
        \[
            L(\theta) = (\# \text{ of red points }\leq \theta) + (\# \text{ of blue points } > \theta)
        \]
        The loss of the $\theta$ shown above is 2, since one red point is to the
        left of $\theta$ and one blue point is to the right.
        Our goal is to design an algorithm for finding a minimizer of
        $L(\theta)$. This is a simple instance of the machine learning task
        of \emph{classification}.

    \begin{subprobset}

        \begin{subprob}
            Also in \python{ep02.py}, write a function named
            \python{compute_ell(data, colors, theta)} which takes in lists
            \python{data} and \python{colors} as described above, as well
            as a floating-point number, \python{theta}. It should return
            the loss at \python{theta} as a floating-point number.  Your
            algorithm should have the best possible time complexity.

            \begin{soln}
                We can write an algorithm that takes time $\Theta(n)$. To compute
                $L(\theta)$, we loop through each of the data points, keeping track of the
                loss incurred. If the data point $x$ is red, but $\leq \theta$, we
                increment the loss by one. If $x$ is blue, but bigger than $\theta$, we
                increment the loss by one. Each check takes constant time and there are
                $\Theta(n)$ points, so the function takes $\Theta(n)$ time in total.

                \inputminted{python}{\thisdir/include-solution/compute_ell.py}

            \end{soln}
            
        \end{subprob}

        \begin{subprob}
            Also in \python{ep02.py}, write a function named \python{minimize_ell(data, colors)}
            which takes in \python{data} and \python{colors} and returns a floating-point
            number which minimizes the loss $L$ for that particular data set.
            Your algorithm should have quadratic time complexity. You may
            assume for simplicity that the smallest data point is
            blue\footnote{Otherwise it is possible (given a special data set) for the
                loss to be minimized at some $\theta$ to the left of all of the
            data.}.

            \begin{soln}
                Note that we can reduce the search space to a finite set ,
                since it suffices to only check the $\theta$ which coincide
                with a data point.  That is, we loop over each of the data
                points $x$ and compute $L(x)$, returning the minimizer.
                Since computing $L(x)$ takes linear time, and there are
                $\Theta(n)$ data points, this approach takes $\Theta(n^2)$
                time in total.

                \inputminted{python}{\thisdir/include-solution/minimize_ell.py}
            \end{soln}
            
        \end{subprob}


        

    \end{subprobset}

\end{progprob}
